% Source: https://artofproblemsolving.com/community/c3h1983256_circle_with_quadrilateral
% Stable identity: aops-topic-1983256-q01
Since $AD$ is a diameter, $\angle ABD=\angle ACD=90^\circ$. Hence $BD=\sqrt{16-1}=\sqrt{15}$ and $AC=\sqrt{16-CD^2}$.

\begin{asy}
usepackage("amsmath");
size(4.5cm);
defaultpen(fontsize(10pt));
pair O=(0,0);draw(circle(O,2));
real al=2*aSin(1/4);
pair A=(-2,0),D=(2,0),B=2*dir(180-al),C=2*dir(180-2*al);
draw(A--B--C--D--cycle,linewidth(1));draw(A--D);draw(B--D,dashed);draw(A--C,dashed);
dot(A);dot(B);dot(C);dot(D);dot(O);
label("$A$",A,W);label("$B$",B,NW);label("$C$",C,N);label("$D$",D,E);label("$O$",O,S);
label("$1$",(A+B)/2,W);label("$1$",(B+C)/2,NW);label("$4$",(1,0),S);
label("\mbox{}",(min(currentpicture,user=true)+max(currentpicture,user=true))/2,black); // keeps the text after the figure black when LaTeX typesets the labels inline
\end{asy}

By Ptolemy's theorem for the cyclic quadrilateral $ABCD$,
\[
AC\cdot BD=AB\cdot CD+BC\cdot AD\quad\Longrightarrow\quad \sqrt{15}\,AC=CD+4 .
\]
Squaring and substituting $AC^2=16-CD^2$: $15(16-CD^2)=(CD+4)^2$, i.e. $16CD^2+8CD-224=0$, or $2CD^2+CD-28=0$, so $(2CD-7)(CD+4)=0$. As $CD>0$, $CD=\boxed{\dfrac{7}{2}}$.
